\chapter{Összefoglalás}
\label{ch:summary}

Szakdolgozatomban a Solid SVG webalapú SVG-szerkesztőt terveztem és valósítottam meg. Az SVG grafikus elemeinek áttekintése és öt elterjedt szerkesztő, az Inkscape, az SVG-Edit, a Boxy SVG, a Figma és a Mediamodifier összehasonlítása alapján hat célt tűztem ki: a szerkesztő telepítés és regisztráció nélkül fusson a böngészőben, kínáljon teljes értékű útvonalszerkesztést logikai műveletekkel, kezelje a többsoros szöveget, jelezze az importált fájlok korlátait, exportáljon szabványos SVG-t és raszteres képet, és legyen telefonon is kezelhető. A célokat tizenöt funkcionális követelményre bontottam, az alkalmazást Svelte keretrendszerrel, TypeScript nyelven készítettem el.

Az elkészült szerkesztő szerver nélkül fut, tizenhárom eszközzel rajzol és szerkeszt, a logikai műveletek eredménye görbékből álló, tovább szerkeszthető alakzat, a szöveg pedig a képernyőn látottnak megfelelően alakítható útvonallá. A tervezés legfontosabb döntései a böngészőre bízott találatvizsgálat, a pillanatképekre épülő előzménykezelés és az egyirányú adatfolyam voltak, ezek tették lehetővé, hogy a működést 447 automatikus teszt ellenőrizze.

A követelmények közül tizenhárom teljesült és kettő részben. Két mérés a legfontosabb célkitűzéseket számszerűen is igazolja. Az útvonallá alakított szöveg képpontjainak legfeljebb 2\%-a tér el a képernyőn látottól, 139 valós SVG-fájl importálásakor pedig a négy ikonkészlet minden fájlja változatlanul került át, és a nagyobb eltérésekről három fájl kivételével figyelmeztetés szólt. A legfontosabb hiányosságok a kódnézet, a minták, vágóútvonalak, maszkok és szűrők kezelése, a szöveg útvonalra illesztése és a vegyes írásirányú szöveg útvonallá alakítása. Ezek, valamint egy magyar nyelvű felület, a továbbfejlesztés irányai.

A munka során a vektorgrafika geometriájától a betűtípusok működéséig sok olyan területet ismertem meg mélyebben, amelyekkel korábban csak felhasználóként találkoztam. A legtöbbet abból tanultam, amikor egy egyszerűnek tűnő funkcionalitás, például a szöveg útvonallá alakítása, a megvalósítás közben bonyolultnak bizonyult.

\chapter*{Summary}
\addcontentsline{toc}{chapter}{Summary}
\label{ch:summary_en}

This thesis presents the design and implementation of Solid SVG, a web-based editor for Scalable Vector Graphics. The work started with an overview of the graphical elements of SVG and a comparison of five widely used editors: Inkscape, SVG-Edit, Boxy SVG, Figma and Mediamodifier. Based on this comparison, six goals were set. The editor should run in the browser without installation, registration or a server; it should offer full path editing with boolean operations that keep curves as curves; it should handle multi-line text and convert it to paths exactly as it appears on screen; it should report which structures of an imported SVG file it cannot handle; it should export standard SVG and raster images; and it should be usable on phones and tablets. These goals were broken down into fifteen functional requirements.

The editor was built with the Svelte framework and TypeScript. Its architecture separates pure geometric functions, the application state and the editing actions, with a one-way data flow from the tools and components to the document. Hit-testing is delegated to the browser, the undo history is based on immutable snapshots, and every change is a single named history step. The thesis describes the more involved parts of the implementation: rotation-aware resizing, three-point circular arcs approximated by cubic Bézier curves, freehand drawing simplified with the Ramer–Douglas–Peucker algorithm, boolean operations and knife cuts with paper.js, stroke-scaled dash patterns, text-to-path conversion with full OpenType shaping, SVG import with explicit warnings, and touch input including pinch zoom.

The behaviour is verified by 447 automated tests, including randomised action sequences, and by manual testing in Chromium-based browsers and Firefox. Two measurements complete the evaluation. Importing 139 real SVG files from seven sources and comparing the rendered original with the rendered export showed that all files of four icon sets survive practically unchanged, and that every larger difference but three is reported to the user; converting text to paths changes at most 2\% of the text's pixels. Thirteen requirements were met in full and two in part. The main missing features are a code view, support for patterns, clipping paths, masks and filters, text on a path, and bidirectional text in path conversion; these, together with a Hungarian user interface, are the directions for further development.

Through this work, I gained a deeper understanding of many areas, from the geometry of vector graphics to how fonts work, that I had previously encountered only as a user. I learned the most when a seemingly simple functionality, such as converting text into a path, turned out to be complicated to implement.